\section{Introduction}
\label{sec:intro}

Image-based disease diagnosis is one of the most-cited use cases of computer vision in agriculture. It is also one of the clearest examples of a benchmark that has outrun deployment.
Models trained on \PV~\cite{hughes2015an}, whose leaves are photographed one at a time on uniform backgrounds, reach about 99\% held-out accuracy~\cite{mohanty2016using}. On photos taken in real fields~\cite{singh2020plantdoc}, the same models do far worse.
Recent work attacks this gap with larger and more specialized models~\cite{salman2025plant,mooraj2026agripath,mamun2025psmamba}. But the numbers being compared are hard to interpret. Label mappings between datasets are rarely published, results usually come from a single run, small field test sets are reported without intervals, and the compute or latency needed in the field is seldom measured.

This paper takes a deliberately modest position. Before proposing new methods for lab-to-field transfer, we should be able to measure it reproducibly, with honest uncertainty, under a compute budget a typical agricultural lab or extension service can afford.
We build such a measurement for the most widely used pair, \PV$\rightarrow$\PD, and use it to ask three practical questions:
\begin{enumerate}
\item How much does generic self-supervised pretraining help compared with small supervised CNNs when \emph{no} field labels exist?
\item How much do a \emph{handful} of field labels (5 per class) buy, compared with all available field labels?
\item Is lab data worth using at all once field labels exist, and does the answer depend on the model?
\end{enumerate}

\paragraph{Contributions.}
\begin{itemize}
\item A released, seeded protocol for \PV$\rightarrow$\PD: a one-to-one 28-class label map, leaf-grouped source splits, a near-duplicate audit (\PDdups of \PDtestN \PD test images duplicate training images), three seeds, bootstrap and paired-bootstrap intervals, and CPU/GPU cost. The full study takes \GPUHours GPU-hours on one consumer GPU (\cref{sec:method,sec:setup}).
\item Five claims pre-registered with numeric thresholds before any results came in. Three hold and two are falsified, and we report both (\cref{sec:results}).
\item Two findings practitioners can act on (\cref{sec:results}). First, a \emph{frozen} DINOv2 ViT-S probe trained only on lab images beats both fine-tuned edge CNNs and fine-tuned DINOv2 on field images. Second, lab pretraining helps CNNs at every label budget we tested. For DINOv2, its value disappears (fine-tuning) or turns negative (lab-trained head for the frozen probe) once all field labels are available.
\item A quantified caution: on \PD, test-set sampling noise ($\pm$\CIhwMin--\CIhwMax\ points) is larger than seed noise, so many published single-number comparisons on this split cannot be distinguished (\cref{sec:limits}).
\end{itemize}
